\section{Related work}
\label{sec:related-work}

The comparison of an implementation with something else that is believed to be right has been studied under four names: benchmarking in model-risk supervision, verification in computational science, the oracle problem in software testing, and agreement between two measurement methods in medical statistics. None of them is written for a risk-model validation engine that is held against printed tables, against closed forms and against a Monte Carlo reference engine at the same time, and each leaves one question of the reconciliation open. This section reads them in that order and ends by placing the proposed protocol among them.

\subsection{Supervisory expectations of benchmarking and code replication}
\label{sec:rw-supervision}

Supervisors ask for the two things that this work turns into measurements. The guidance on model risk management of the Federal Reserve and the Office of the Comptroller of the Currency defines benchmarking as the comparison of a model's inputs and outputs with estimates from alternative internal or external data or models. It states that the benchmark is itself an alternative prediction whose differences from the model may come from different data or methods, and that when the model and the benchmark match well this is evidence in favour of the model, to be interpreted with caution so that the bank does not get a false degree of comfort \citep{fed2011sr117}. The guide to internal models of the European Central Bank lists, among the minimum analyses of the validation of a credit rating system, a quality assurance of the computer code that shows the implementation to be compliant with, and to reproduce exactly, the documented model under review \citep{ecb2025guide}. The working paper of the Basel Committee on the validation of internal rating systems is about the calibration and discrimination tests themselves, which the engine implements, and not about the check of an implementation of them \citep{bcbswp14}.

Three gaps follow from the passages cited. They do not say what a match is: how close, by which statistic, on how many points. They do not say what a match is evidence of when the benchmark shares code, parameters or the same library with the model. And they do not say what the reader is told when the match fails. The protocol of Section~\ref{sec:proposed-method} answers each of the three with a rule that is stated before the comparison and a statement of what the comparison does not show.

\subsection{Verification, validation and the test oracle}
\label{sec:rw-verification}

Computational science separates verification, which asks whether the equations are solved correctly, from validation, which asks whether the right equations are solved. Roache treats the quantification of uncertainty in computational fluid dynamics \citep{roache1997}, and Oberkampf and Trucano review verification and validation in the same field \citep{oberkampf2002}. Verification of a code, as distinct from verification of one calculation, evaluates the error of the code against known benchmark solutions, and the method of manufactured solutions is a general procedure for generating analytical solutions to serve as such benchmarks \citep{roache2002}. The reconciliations of this work are verifications in this sense: they ask whether the engine computes the number that a printed table, a second implementation or a closed form says it should, and say nothing on whether Hull--White or the Jeffreys test is the right model for a bank.

Software testing calls the same difficulty the oracle problem, the absence of a mechanism that tells the correct output of a program for an arbitrary input, and the literature on it is surveyed by Barr and co-authors \citep{barr2015}. Metamorphic testing replaces the oracle by metamorphic relations, necessary properties that the outputs of related inputs must satisfy \citep{chen2018}. Differential testing runs independent implementations on the same input and studies the disagreements; it found previously unknown bugs in C compilers by compiling randomly generated programs \citep{yang2011}. A printed table is a specified oracle with a rounding, a second implementation a differential oracle, and ORE or QuantLib a differential oracle with its own Monte Carlo error or its own numerical tolerance. The literature on oracles does not treat the case in which the oracle and the program call the same library, which is the first of the cases measured in Section~\ref{sec:results}.

\subsection{Agreement between two measurements}
\label{sec:rw-agreement}

Medical statistics has a long literature on whether two methods measure the same quantity. Bland and Altman argued that the correlation coefficient is the wrong tool and proposed the limits of agreement, the interval that contains most differences between the two methods \citep{bland1986}; Lin proposed a concordance correlation coefficient to evaluate reproducibility, which penalises location and scale shifts as well as scatter \citep{lin1989}. Both describe the differences of paired measurements; neither fixes an acceptance bound, because the bound belongs to the application. A reconciliation with a reference engine has the same structure with one difference: the two sides are not measurements of an unknown but computations with a known sampling error, so the difference can be standardised before it is judged.

Judging a difference is a test, and the test must be chosen for its question. A failure to reject a difference does not establish equality, a point made in a statistics note by Altman and Bland \citep{altman1995}; equivalence is established by two one-sided tests against a stated margin \citep{schuirmann1987}, a procedure that Lakens presents in a practical primer \citep{lakens2017}. The acceptance rules of risk-model reconciliations, including the rule of the product studied here, are difference tests ``within two standard errors''. Section~\ref{sec:results} reads the committed comparisons both ways and reports the equivalence margin that each supports.

\subsection{Monte Carlo error and simultaneous comparison}
\label{sec:rw-montecarlo}

The comparison of two Monte Carlo estimates needs the standard error of each. Glasserman gives the estimators and the variance reduction by antithetic pairs \citep{glasserman2004}; Koehler, Brown and Haneuse treat the assessment of Monte Carlo error in simulation-based analyses \citep{koehler2009}, and Morris, White and Crowther give the Monte Carlo standard error of the performance measures of a simulation study \citep{morris2019}. A rate measured by repetitions carries a Wilson interval \citep{brown2001}, which this work uses for every rejection rate.

A reconciliation compares a profile at many dates and several scalar functionals at once. The probability that all the comparisons pass when the engines agree falls with their number, and the bound that restores the family-wise size for independent comparisons is the one Sidak proved for rectangular confidence regions of multivariate normal means; his inequality holds for any correlation of the comparisons \citep{sidak1967}. The reconciliation rules in use fix the bound per comparison; Section~\ref{sec:theory} derives what that does to the size of the family.

\subsection{The models, engines and papers reconciled}
\label{sec:rw-models}

The exposure engine implements the one-factor Hull--White model of Hull and White \citep{hull1990}, whose numerical treatment they describe in a later paper \citep{hull1994}; the exact decomposition of an option on a coupon bond into options on zero-coupon bonds is due to Jamshidian \citep{jamshidian1989}, and Brigo and Mercurio treat the model in book form \citep{brigo2006}. The expected exposure, the effective expected positive exposure and the credit valuation adjustment follow the treatments of Gregory \citep{gregory2020} and of Brigo, Morini and Pallavicini \citep{brigo2013}; the margin period of risk is modelled after Andersen, Pykhtin and Sokol \citep{andersen2017}; and the wrong-way risk calibration is the hazard-rate model of Hull and White \citep{hullwhite2012}. The reference engines are the Open Source Risk Engine \citep{ore2026} and QuantLib \citep{quantlib2025}; a book-length treatment of derivatives pricing and credit exposure analysis is that of Lichters, Stamm and Gallagher \citep{lichters2015}. The printed answers belong to five papers: the most prudent estimate of a default probability in a low-default portfolio \citep{plutotasche2005}, the calibration of a default-probability curve to a new year \citep{tasche2013}, the generator of a Markov chain from an empirical transition matrix \citep{israel2001}, the interval estimates of default probabilities \citep{schuermann2004}, and the stress-testing projection of default probabilities from a transition matrix \citep{engelmann2024}.

Two software libraries underlie every number: SciPy \citep{virtanen2020} and NumPy \citep{harris2020}. Floating-point arithmetic limits what a comparison can resolve \citep{goldberg1991,higham2002}, and Gauss quadrature, which the independent route for the Jeffreys test uses, is compared with Clenshaw--Curtis quadrature in its convergence by Trefethen \citep{trefethen2008}. The reproducibility of such work is the subject of the rules for reproducible computational research \citep{peng2011,sandve2013}.

\subsection{Where the proposed protocol stands}
\label{sec:rw-position}

The literature above supplies the parts and no whole. Supervisors say that a benchmark must be used with caution and do not define a match. Verification supplies benchmark solutions for equation solvers, and software testing supplies the oracle problem and differential testing. Medical statistics supplies agreement and equivalence statistics, and the simulation literature supplies the error of a Monte Carlo estimate and the cost of simultaneous comparisons. The protocol proposed here fixes, for a validation engine, one rule per kind of reference, the statistic and bound of each, the listing of every gap with its explanation, and the list of what is not claimed; and it measures the size and power of the acceptance rule that the product uses for the Monte Carlo reference, which none of the cited works does for a risk engine. The protocol is not a new statistic: it is a combination of known ones into a procedure that a validation function can run, audit and report.
